\documentclass[11pt]{article}

\usepackage[margin=25mm]{geometry}

\usepackage{amsmath,amssymb,fontspec,unicode-math}

\IfFontExistsTF{DejaVu Serif}{\setmainfont{DejaVu Serif}}{\setmainfont{Latin Modern Roman}}

\IfFontExistsTF{Latin Modern Math}{\setmathfont{Latin Modern Math}}{\setmathfont{DejaVu Math TeX Gyre}}

\usepackage{xurl,hyperref}

\setlength{\parindent}{0pt}\setlength{\parskip}{6pt}

\begin{document}

\title{MAPLE: Multi-scale Attribute-enhanced Prompt Learning for Few-shot Whole Slide Image Classification}

\author{ELVIA preliminary AI analysis}\date{21 October 2026}\maketitle

Preliminary assessment: Minor.\par

Independent scientific review remains pending. Author adoption has not been requested.

The release pairs entity-major prompts with a class-major reshape and implements training/graph choices that differ from the displayed method. Exact feature witnesses establish those local differences; reported classification performance and patient splits are not invalidated.

\section*{Review scope}

Claim-focused reading of the original paper and relevant author-released passages; direct mathematical or static-source reasoning. No author program, earlier verifier, model or benchmark was executed during this check.

\section*{Problem 1: Entity-major prompts are consumed as class-major prompts}

Source: Eq.(7), PDF p.6; model.py, lines 58–78,479–490

Assessment: supported; impact: minor

\subsection*{Source claim (paraphrase)}

Prompt construction uses flattened index \(eC+c\), while the reshape/permutation reads \(cE+e\) for entity \(e\), class \(c\).

\subsection*{Analysis}

For \(E=8,C=2\), only the first and final semantic slots retain their indices, leaving \(14\) of \(16\) mismatched; \(C=3\) gives \(22\) of \(24\). With unit zero-mean vector \(u\in\mathbb R^{512}\), visual features all \(u\), class prompts \(-u,u\), and unit logit scale, proper pairing gives mean logits \((-1,1)\), probabilities approximately \((0.119203,0.880797)\). The release groups alternating prompt rows into both class columns, whose means are \((0,0)\), probabilities \((0.5,0.5)\).

\subsection*{Finding and scope}

The feature-space computation loses semantic pairing and can change predictions. Actual-slide embedding realizability, adaptation to old wiring and benchmark effect remain separate.

\subsection*{Proposed correction}

Reshape the entity-major array directly to entity, class, embedding dimensions, without the class-major reshape/permutation. This restores the index identity under the stated construction.

A proposed correction is retained where stated; independent verification and author adoption remain unestablished.

\section*{Problem 2: Stop-gradient and zero-fusion loss change the stated objective}

Source: Eqs.(15)–(16), PDF p.7; model.py, lines 487–488,559–568

Assessment: supported; impact: minor

\subsection*{Source claim (paraphrase)}

The paper applies cross entropy to fused entity/slide logits; the release detaches visual entity features for positive fusion weight and changes the loss at zero weight.

\subsection*{Analysis}

After holding semantic indexing equal, take entity logits \((-\sin\theta,\sin\theta)\), constant zero slide logits, target class \(1\), and \(\lambda=0.3\). The displayed loss derivative at \(\theta=0\) is \(-(1-\lambda)=-0.7\), while the detached path gives \(0\). Attributes and the separate slide branch can still learn. At \(\lambda=0\), half the entities with logits \((1,-1)\), half with \((-1,1)\), give fused CE \(\log2\approx0.693147\), but average individual CE \(\log(2\cosh1)\approx1.126928\). The latter may be a deliberate stronger supervision choice, but is a different objective.

\subsection*{Finding and scope}

The release does not differentiate the literal stated loss in those branches. Neither choice is proved inferior or a cause of the reported scores.

\subsection*{Proposed correction}

Either document the stop operator and separate endpoint loss, or follow ordinary fused-logit cross entropy throughout. No accuracy recovery is claimed.

A proposed correction is retained where stated; independent verification and author adoption remain unestablished.

\section*{Problem 3: The released graph is text-fixed rather than slide-dependent}

Source: Eq.(8), PDF p.6; Figure 2, PDF p.4; model.py, lines 102–128,495–523,572

Assessment: supported; impact: minor

\subsection*{Source claim (paraphrase)}

The equation builds neighbors from slide visual entity features, while the release saves generic-text adjacency once.

\subsection*{Analysis}

Take \(16\) nodes, first eight generic-text vectors \(u\), remaining eight orthogonal vector \(v\), with seven nonself neighbors plus a self-loop. Both text groups have an unambiguous similarity boundary. A slide initially matching text has the same adjacency. Swapping visual vectors of nodes \(0\) and \(8\) makes the visual-neighbor set for node \(0\) equal to \(\{0,9,\ldots,15\}\), but the release still uses \(\{0,\ldots,7\}\). Its current visual inputs set node features, not edges. Figure 2 partly supports a text-graph interpretation, so author intent is ambiguous.

\subsection*{Finding and scope}

The graph specification and release differ; fixed text relations may still be useful, and graph attention is present.

\subsection*{Proposed correction}

Specify a fixed semantic-text graph with explicit conventions, or rebuild the declared visual graph per slide. The evidence does not select the better empirical variant.

A proposed correction is retained where stated; independent verification and author adoption remain unestablished.

\section*{Limits}

Preliminary AI assessment; independent human scientific adjudication remains pending.\par

Analytic witnesses and static release checks do not remeasure the published benchmark results.\par

\section*{Source identity}

NeurIPS 2025 official proceedings PDF (retained local copy)

\url{https://proceedings.neurips.cc/paper_files/paper/2025/file/a94a8800a4b0af45600bab91164849df-Paper-Conference.pdf}

Reviewed PDF SHA-256: \nolinkurl{32df3643d1f888fde9751647eb282f56112fccf0fe5962eb8c9aae814ff8c431}

The SHA-256 identifies the exact retained PDF inspected in this rapid source/argument check. Original supplied report bytes are retained locally. No unambiguous original PDF retrieval timestamp was recovered; file modification time was not substituted.

\end{document}
